\documentclass[conference]{IEEEtran}
\IEEEoverridecommandlockouts
% The preceding line is only needed to identify funding in the first footnote. If that is unneeded, please comment it out.
\usepackage{cite}
\usepackage{amsmath,amssymb,amsfonts}
\usepackage{algorithmic}
\usepackage{graphicx}
\usepackage{textcomp}
\usepackage{xcolor}
\usepackage{booktabs}
\usepackage{tcolorbox}
\tcbuselibrary{skins}
\usepackage{tabularx}
\def\BibTeX{{\rm B\kern-.05em{\sc i\kern-.025em b}\kern-.08em
    T\kern-.1667em\lower.7ex\hbox{E}\kern-.125emX}}
\begin{document}

\title{An End-to-End Multimodal AI System for Alzheimer's Disease Staging: Harmonized GNN Classification and Knowledge-Grounded Report Generation}

\author{
\IEEEauthorblockN{Wei-Chi Lin}
\IEEEauthorblockA{\textit{Miin Wu School of Computing} \\
\textit{National Cheng Kung University}\\
Tainan, Taiwan \\
nm6131027@gs.ncku.edu.tw}
\and
\IEEEauthorblockN{Ya-Ning Chang}
\IEEEauthorblockA{\textit{Miin Wu School of Computing} \\
\textit{National Cheng Kung University}\\
Tainan, Taiwan \\
yaningchang@gs.ncku.edu.tw}
}

\maketitle

\begin{abstract}
The early and accurate staging of Alzheimer's Disease (AD) is crucial for timely clinical intervention. While neuroimaging modalities like fMRI and sMRI provide complementary insights, transitioning multimodal AI models into clinical workflows is hindered by multi-site data confounding, a lack of functional-network interpretability, and the absence of fact-grounded diagnostic reports. In this paper, we propose a comprehensive end-to-end multimodal AI pipeline. We demonstrate that label-free ComBat harmonization successfully mitigates severe site-class imbalance, revealing true disease signals in functional networks. Our multimodal architecture leverages a frozen BrainIAC Vision Transformer for sMRI and a Graph Attention Network (GAT) with virtual node pooling for fMRI. Features are fused via a Patient-Specific Cross-Attention Generation (PCAG) mechanism with per-patient sigmoid gating. Finally, we implement a Knowledge-Grounded Retrieval-Augmented Generation (RAG) system utilizing a Neo4j Knowledge Graph to produce diagnostic reports grounded in explicit model evidence and retrieved medical literature. Evaluations demonstrate significant improvements in clinical relevance and completeness of generated reports.
\end{abstract}

\begin{IEEEkeywords}
Alzheimer's Disease, Multimodal Data Fusion, Graph Neural Networks, Site Harmonization, Retrieval-Augmented Generation, Knowledge Graph
\end{IEEEkeywords}

\section{Introduction}

The early and accurate staging of Alzheimer's Disease (AD)---distinguishing among Normal Control (NC), Mild Cognitive Impairment (MCI), and AD---is critical for timely clinical intervention. Neuroimaging modalities, particularly functional MRI (fMRI) and structural MRI (sMRI), provide complementary insights into functional connectivity disruptions and anatomical atrophy, respectively. However, transitioning multimodal AI models into clinical workflows remains obstructed by three bottlenecks: multi-site data confounding, a lack of network-level functional interpretability, and the absence of fact-grounded clinical diagnostic reports.

Large-scale neuroimaging cohorts are frequently assembled across multiple clinical sites, introducing severe site-class imbalances. In combined cohorts such as ADNI and TPMIC, one site may be predominantly populated by NC cases while another is skewed toward disease stages. Under such imbalances, models inadvertently learn ``site signatures'' (e.g., scanner-specific artifacts) rather than pathological markers, achieving artificially high cross-site performance while catastrophically failing within-site evaluations. Standard label-aware harmonization techniques fail in these scenarios, preserving site-related variation under the guise of disease signal.

Despite achieving high predictive AUC scores, state-of-the-art architectures commonly function as black boxes. For a staging system to be clinically viable, physicians require explicit functional-network or structural evidence rather than abstract latent representations. Furthermore, while LLMs demonstrate immense potential for clinical report generation, they are prone to factual hallucination \cite{ji2023hallucination} and produce narratives disconnected from concrete model evidence or verified medical literature.

To address these gaps, we propose an end-to-end multimodal AI system for AD staging, integrating label-free site harmonization, an interpretable multimodal fusion architecture, and a Knowledge-Grounded RAG system for automated reporting.

The main contributions are:
\begin{itemize}
    \item \textbf{Label-Free Site Harmonization:} We show that under severe site-class imbalance, label-aware ComBat inadvertently preserves site signatures. Label-free ComBat on fMRI functional connectivity eliminates pure site effects while preserving disease signals, evidenced by improved within-site AUC and a shift toward disease-dominant signal in the Default Mode (DMN), Sensorimotor (SMN), Salience (SN), and Cerebellar (CereN) networks after harmonization.
    \item \textbf{Multimodal Fusion with Network Interpretability:} We adapt PCAG \cite{zhang2024pcag} to multimodal neuroimaging, pairing a GAT encoder with 9-network pooling for fMRI---exposing resting-state network-level functional evidence---with a frozen BrainIAC ViT for sMRI. Setting fMRI as the Query encodes the neuroimaging prior that functional connectivity state determines which structural features are salient. Predictions are further stabilized by an OOF variance-penalized ensemble strategy across five random seeds.
    \item \textbf{Knowledge-Grounded Report Generation:} We implement a Neo4j Knowledge Graph paired with MMR retrieval, combining retrieved medical literature with functional network evidence to prompt Gemma3-12B for fact-grounded clinical reports.
    \item \textbf{Multi-Layer Evaluation Framework:} We establish a rigorous validation protocol for LLM-generated reports, combining an LLM judge, classification-consistency validity checks, and objective citation metrics to quantify RAG-driven improvements in clinical relevance and completeness.
\end{itemize}

\section{Related Work}

\subsection{Deep Learning for Alzheimer's Disease Classification}
The application of deep learning to neuroimaging has significantly advanced AD staging. Early architectures such as BrainNetCNN \cite{kawahara2017brainnetcnn} pioneered CNNs for topological brain connectivity matrices, while GNN-based models including BrainGNN \cite{li2021braingnn} and Hi-GCN \cite{jiang2020higcn} established strong baselines on functional connectomes. More recently, meta-learning frameworks (e.g., ADMGCN \cite{sun2025admgcn}) and spatio-temporal GNNs modeling dynamic functional connectivity \cite{wen2025stdfbn} have further advanced multi-class AD staging. Despite their accuracy, a critical limitation of these models is their reliance on small, single-site cohorts (e.g., $n \leq 85$), leaving cross-site generalizability largely unaddressed.

\subsection{Multi-Site Data Harmonization and Confounding}
Aggregation of neuroimaging data across sites introduces non-biological scanner biases. ComBat---originally developed for genomics batch correction \cite{johnson2007combat}---has become the standard for harmonizing structural and functional MRI \cite{fortin2018harmonization}. CovBat \cite{chen2022covbat} extends ComBat by additionally harmonizing covariance structure via principal component decomposition. However, both methods assume approximately balanced site-class distributions. When severe imbalance exists, label-aware harmonization preserves site-correlated signals as biological variance, creating classification shortcuts that inflate cross-site performance while failing within-site---necessitating label-free alternatives.

\subsection{Interpretable AI in Medical Imaging}
While deep learning achieves high classification metrics, its black-box nature limits clinical adoption. GradCAM \cite{selvaraju2017grad} highlights spatially localized diagnostic regions in structural imaging. For functional graph models, layer-wise attention maps in GATs expose network-level evidence. For Transformer-based encoders, Attention Rollout \cite{abnar2020rollout} propagates attention across all layers, providing more faithful explanations than single-layer scores. The persistent challenge is unifying these modality-specific interpretations into coherent, actionable clinical insights.

\subsection{Knowledge-Grounded RAG for Clinical Reporting}
Retrieval-Augmented Generation (RAG) \cite{lewis2020rag} grounds LLM outputs in retrieved documents, mitigating hallucination. Knowledge Graph (KG)-enhanced retrieval has been shown to outperform flat vector databases in medical settings: MedRAG \cite{xiong2025medrag} integrates hierarchical diagnostic KGs for clinical reasoning, MedGraphRAG \cite{wu2025medgraphrag} builds triple-based graphs linking patient records to curated literature for evidence-grounded medical QA, and LAB-KG \cite{labkg2025} conditions generation on KG-retrieved patient conditions for lab test interpretation. Our system extends this paradigm by grounding KG-RAG reporting directly in functional network-level evidence from an upstream multimodal GNN, creating a coherent chain from neuroimaging to diagnostic narrative.

\section{Methodology}

Fig.~\ref{fig:arch} illustrates the complete pipeline. We detail each component in the following subsections.

% FIGURE 1: System Architecture
% Upload as: architecture.png (or architecture.pdf)
\begin{figure}[t]
  \centering
  \includegraphics[width=\columnwidth]{architecture}
  \caption{Overview of the proposed end-to-end pipeline. fMRI functional connectivity graphs and sMRI structural scans are encoded by a GAT and a frozen BrainIAC ViT, respectively. The PCAG mechanism fuses both modalities via patient-specific sigmoid-gated cross-attention. The predicted class and GAT attention evidence are used to query a Neo4j Knowledge Graph, and MMR retrieval of relevant literature together prompts Gemma3-12B to generate a fact-grounded clinical report.}
  \label{fig:arch}
\end{figure}


\subsection{Dataset and Multi-Site Harmonization}
The pipeline is evaluated on an aggregated cohort of ADNI \cite{petersen2010adni} and TPMIC, totaling 204 patients (155 training / 49 held-out test; NC=87, MCI=72, AD=45). The two sites exhibit severe class imbalance: in the training set, ADNI subjects are 70\% NC (NC=39, MCI=14, AD=3), while TPMIC subjects are 69\% disease-stage (NC=31, MCI=55, AD=13), making site identity a strong confound for disease classification.

Standard label-aware ComBat \cite{fortin2018harmonization} fails under such imbalance, inadvertently preserving site-specific signatures as biological variance. We instead apply \textit{label-free ComBat} globally on the 6670-dimensional upper triangle of each patient's fMRI functional connectivity (FC) matrix, computed across 116 AAL \cite{tzourio2002aal} ROIs. This removes pure site effects without conditioning on diagnosis, thereby preventing site-class confounding. For sMRI, label-aware ComBat is applied per fold to prevent data leakage while retaining within-fold disease signal.

\subsection{Multimodal Encoders}
\textbf{fMRI Graph Encoder:} Each patient's brain is modeled as a graph of 116 AAL ROI nodes. Each node's input feature is the corresponding row of the FC matrix (116-dimensional), encoding its pairwise connectivity with all other ROIs. The graph is sparsified at $K_{\text{ratio}} = 0.20$, retaining the top 20\% of edges by connectivity strength. A 3-layer, 4-head GAT \cite{velickovic2018gat} with hidden dimension 128 processes the graph; a virtual node---connected to all ROI nodes---aggregates global context into a 128-dimensional readout. Node embeddings are further pooled into nine canonical resting-state networks: Default Mode (DMN), Sensorimotor (SMN), Visual (VN), Salience (SN), Frontoparietal (FPN), Language (LN), Ventral Attention (VAN), Basal Ganglia (BGN), and Cerebellar (CereN). Concatenating the nine network embeddings with the virtual node embedding yields a 1280-dimensional functional representation. Layer-wise GAT attention maps expose network-level functional evidence; following label-free harmonization, DMN, SMN, SN, and CereN exhibit disease-dominant signal over site artifacts.

\textbf{sMRI Structural Encoder:} Structural features are extracted via BrainIAC \cite{tak2026brainiac}, a contrastive self-supervised ViT-B pretrained on 32,015 multiparametric brain MRI scans across 10 medical conditions. BrainIAC is used frozen, producing a 768-dimensional structural embedding per patient.

\subsection{Patient-Specific Cross-Attention Fusion (PCAG)}
\label{sec:pcag}
We adapt the Patient-Specific Cross-Attention Generation (PCAG) module \cite{zhang2024pcag} for multimodal neuroimaging fusion. PCAG addresses a known weakness of standard cross-attention \cite{vaswani2017attention}: when cross-modal interactions are uninformative, conventional attention produces noisy representations that degrade downstream task performance. PCAG mitigates this through dual sigmoid gating---a pre-gate that suppresses non-informative interactions and a contextual gate that reduces residual uncertainty.

In our adaptation, we treat the 1280-dimensional fMRI embedding as the Query and the 768-dimensional sMRI embedding as Key and Value. Linear projections $W_Q \in \mathbb{R}^{d_f \times 1280}$ and $W_K, W_V \in \mathbb{R}^{d_f \times 768}$ (with $d_f = 20$) map both modalities into a shared fusion space. The patient-specific gating is:
\begin{equation}
\mathbf{S} = \sigma\!\left(\frac{(W_Q \mathbf{e}_f) \odot (W_K \mathbf{e}_s)}{\sqrt{d_f}}\right), \quad \mathbf{h}_{\text{fused}} = \mathbf{S} \odot (W_V \mathbf{e}_s)
\end{equation}
where $\sigma(\cdot)$ is element-wise sigmoid and $\odot$ is element-wise multiplication. The gating vector $\mathbf{S} \in \mathbb{R}^{d_f}$ assigns per-patient, per-dimension weights that adaptively modulate the structural modality's contribution. Setting fMRI as the Query reflects a neuroimaging prior: a patient's functional connectivity state determines which structural features are clinically salient---an asymmetry absent from symmetric attention. The compact $d_f=20$ fusion space further forces PCAG to capture fine-grained inter-modal dependencies that concatenation-based fusion resolves only globally. Fig.~\ref{fig:pcag} illustrates the complete data flow.

% FIGURE: PCAG mechanism diagram
% Upload as: pcag_diagram.png
\begin{figure}[t]
  \centering
  \includegraphics[width=\columnwidth]{pcag_diagram}
  \caption{PCAG fusion mechanism (see §\ref{sec:pcag}). The fMRI embedding is projected to a Query vector; the sMRI embedding is projected to Key and Value vectors. An element-wise sigmoid gate $\mathbf{S}$ computed from the scaled Query--Key product is applied to the Value, yielding a 20-dimensional per-patient fused representation.}
  \label{fig:pcag}
\end{figure}

\subsection{Ensemble Training and Variance-Penalized Selection}
Models are trained with 5-fold cross-validation over 200 epochs, optimized by AdamW (lr=$3\times10^{-4}$) with a class-balanced sampler. For the NC vs.\ AD task only, FC-Mixup and DropEdge augmentation is applied; augmentation is omitted for NC vs.\ MCI and MCI vs.\ AD, where inter-class boundaries are subtle and mixing-based augmentation was empirically found to degrade performance.

Small held-out test sets ($n{=}49$) produce high-variance AUC estimates: a single poorly-seeded model can swing the test AUC by $\pm$0.05--0.10, rendering single-run comparisons unreliable. To resist this evaluation noise, we train five seeds (42, 123, 456, 789, 2024) and select the aggregation policy per task using an OOF variance-penalized criterion \cite{caruana2004ensemble} that explicitly penalizes strategies prone to seed-selection instability:
\begin{equation}
\text{Score} = \text{OOF\_AUC} - \lambda \cdot \sigma_{\text{seed}} \cdot \alpha
\end{equation}
where $\lambda = 0.5$, $\sigma_{\text{seed}}$ is the standard deviation of per-seed OOF AUC, and $\alpha \in \{1.0, 0.5, 1.5\}$ for mean, median, and top-3 aggregation respectively. The penalty discourages high-variance strategies, selecting median aggregation for MCI tasks and top-3 seeds for NC vs.\ AD.

\subsection{Knowledge-Grounded RAG System}
The RAG pipeline has two offline components. \textbf{Patient Knowledge Graph:} All training patients' GNN embeddings are stored as nodes in Neo4j, with \textit{SIMILAR\_TO} edges formed by cosine similarity, encoding population-level functional-structural similarity patterns. \textbf{Literature Index:} AD-relevant medical literature is embedded via \textit{nomic-embed-text} and indexed as a Neo4jVector collection within the same database.

At inference, the patient's predicted class and extracted functional network indicators are combined into a semantic query. Maximal Marginal Relevance (MMR) retrieval \cite{lewis2020rag} selects $k=5$ diverse documents from 20 candidates, balancing relevance and coverage. The retrieved literature together with the patient's network-level attention evidence is structured into a prompt for Gemma3-12B \cite{gemma3_2025} (deployed via Ollama), which generates a fact-grounded clinical diagnostic report.

\section{Results}

% FIGURE 2: ROC Curves
% File to upload: roc_curves_v2_nolabel.png -> rename to roc_curves.png
\begin{figure}[t]
  \centering
  \includegraphics[width=\columnwidth]{roc_curves}
  \caption{ROC curves for all three binary classification tasks on the held-out test set ($n{=}49$). Shaded areas indicate 95\% bootstrap confidence intervals.}
  \label{fig:roc}
\end{figure}

% FIGURE 3: Site Confound Analysis
% File to upload: site_confound_analysis_v2.png -> rename to site_confound.png
\begin{figure}[t]
  \centering
  \includegraphics[width=\columnwidth]{site_confound}
  \caption{Within-site AUC before and after label-free ComBat harmonization across all three tasks. NC vs.\ MCI (ADNI) rises from 0.222 to 0.444 ($\Delta{=}{+}0.222$), confirming that harmonization eliminates site-identity classification and restores genuine disease-signal discrimination.}
  \label{fig:siteconfound}
\end{figure}

% FIGURE 4: Report Quality
% File to upload: report_quality_v5_full_fix.png -> rename to report_quality.png
\begin{figure}[t]
  \centering
  \includegraphics[width=\columnwidth]{report_quality}
  \caption{LLM-judge evaluation of RAG vs.\ no-RAG reports ($n{=}49$). Panels (A–B): Gemma3 mean scores and $\Delta$ with 95\% CI. Panels (C–D): Qwen3. RAG significantly improves clinical relevance (Qwen3: $\Delta{=}{+}0.333$, $p{=}0.007$). No significant coherence penalty is observed ($p{>}0.17$ for both judges).}
  \label{fig:reportquality}
\end{figure}

\subsection{Classification Performance}

Table~\ref{tab:auc} and Fig.~\ref{fig:roc} report held-out test AUC with 95\% bootstrap confidence intervals for the OOF-selected 5-seed ensemble. Our system achieves AUC of 0.791, 0.686, and 0.672 for NC vs.\ AD, NC vs.\ MCI, and MCI vs.\ AD, respectively. NC vs.\ MCI is universally the most challenging task, reflected in the widest confidence intervals. We note that overall NC vs.\ MCI AUC is lower than the pre-harmonization baseline (0.747); this decrease is intentional and expected---as demonstrated in §\ref{sec:site}, the baseline inflates cross-site performance by exploiting scanner identity rather than disease signal. The harmonized system's lower but clinically honest AUC reflects genuine disease-signal discrimination.

\begin{table}[h]
\caption{Ensemble Classification Performance (test set, $n{=}49$)}
\label{tab:auc}
\centering
\begin{tabular}{lcc}
\toprule
Task & AUC & 95\% CI \\
\midrule
NC vs AD  & 0.791 & [0.609, 0.927] \\
NC vs MCI & 0.686 & [0.509, 0.855] \\
MCI vs AD & 0.672 & [0.468, 0.879] \\
\bottomrule
\end{tabular}
\end{table}

We note that truly standardized multi-site benchmarks for fMRI+sMRI AD classification remain scarce in the literature: most existing methods either evaluate on single-site ADNI or use additional modalities (DTI, PET) unavailable in our cohort. Table~\ref{tab:benchmark} therefore contextualizes our system against the three most representative configurations available, spanning single-site and multi-site setups. ADMGCN~\cite{sun2025admgcn} is the most directly comparable baseline (same ADNI fMRI, GCN, all three tasks): we outperform it on NC vs.\ AD (+0.032) and NC vs.\ MCI (+0.021). Yuan et al.~\cite{yuan2026multimodal} applies ComBat harmonization to a similarly-sized multi-site cohort ($n{=}234$) but includes an additional DTI modality; our system achieves competitive performance on NC vs.\ AD and NC vs.\ MCI without DTI, demonstrating that effective harmonization can compensate for the absence of white-matter structural connectivity features. Wen et al.~\cite{wen2025stdfbn} reports the highest AUC on a small single-site cohort ($n{=}85$), but the absence of multi-site validation limits its generalizability. Our MCI vs.\ AD AUC (0.672) is weakest across methods; given the small AD test subset ($n{=}18$, CI [0.468, 0.879]), this gap should be interpreted with caution.

\begin{table}[h]
\caption{Comparison with Related Methods (AUC)}
\label{tab:benchmark}
\centering
\begin{tabular}{lcccc}
\toprule
Method & Sites & NC/AD & NC/MCI & MCI/AD \\
\midrule
\textbf{Ours (PCAG-ComBat)}        & 2 (ADNI+TPMIC) & \textbf{0.791} & \textbf{0.686} & 0.672 \\
ADMGCN \cite{sun2025admgcn}        & 1 (ADNI)       & 0.759 & 0.665 & \textbf{0.758} \\
Yuan et al.~\cite{yuan2026multimodal}$^{\dag}$ & $\geq$2        & 0.901 & 0.839 & 0.809 \\
Wen et al.~\cite{wen2025stdfbn}$^{\ddag}$     & 1 (ADNI)       & 0.831 & 0.792 & 0.769 \\
\bottomrule
\multicolumn{5}{l}{\small $^{\dag}$Uses additional DTI modality ($n{=}234$). $^{\ddag}$Small single-site cohort ($n{=}85$).}
\end{tabular}
\end{table}

\subsection{Ablation Study}

Table~\ref{tab:ablation} presents single-seed (seed=42) AUC for ablation variants; all share the same train/test split for fair comparison.

\begin{table}[h]
\caption{Ablation Study (seed=42, $n{=}49$ test)}
\label{tab:ablation}
\centering
\begin{tabular}{lccc}
\toprule
Method                               & NC/AD          & NC/MCI         & MCI/AD \\
\midrule
\multicolumn{4}{l}{\textit{Single-seed ablation (seed=42, for component comparison)}} \\
SVM (FC features)                    & 0.755          & 0.750          & 0.505 \\
Unimodal (fMRI only)                 & 0.768          & 0.683          & 0.646 \\
Unimodal (sMRI only)                 & 0.750          & 0.533          & 0.515 \\
Concatenation Fusion                 & \textbf{0.804} & 0.417          & 0.535 \\
Without Harmonization                & 0.682          & 0.678          & 0.525 \\
Inverted Cross-Attention (Q=sMRI)    & 0.700          & 0.681          & 0.515 \\
\textbf{PCAG-ComBat (Proposed)}      & 0.682          & \textbf{0.708} & \textbf{0.651} \\
\midrule
\multicolumn{4}{l}{\textit{5-seed ensemble comparison (same protocol as final results)}} \\
fMRI-only ensemble                   & 0.777          & 0.675          & 0.500 \\
\textbf{PCAG-ComBat ensemble}        & \textbf{0.791} & \textbf{0.686} & \textbf{0.672} \\
\bottomrule
\end{tabular}
\end{table}

\textbf{Modality contribution.} Both fMRI and sMRI contribute complementary information: fMRI-only outperforms sMRI-only on MCI vs.\ AD (0.646 vs.\ 0.515), reflecting the utility of functional connectivity for subtle staging, while sMRI-only remains competitive on NC vs.\ AD.

\textbf{Harmonization effect.} Removing ComBat substantially degrades both MCI tasks (NC vs.\ MCI: $0.708{\to}0.678$; MCI vs.\ AD: $0.651{\to}0.525$), confirming that harmonization is critical for learning genuine disease signal rather than site identity.

\textbf{Fusion design.} PCAG achieves the highest AUC on NC vs.\ MCI and MCI vs.\ AD, where subtle inter-class differences demand patient-specific cross-modal integration. Concatenation Fusion leads on NC vs.\ AD (0.804 vs.\ 0.682 for single seed), but this single-seed comparison understates PCAG's capability: the full 5-seed ensemble recovers NC vs.\ AD to 0.791 (Table~\ref{tab:auc}). Furthermore, Concatenation Fusion collapses to below-chance on NC vs.\ MCI (0.417), demonstrating that it cannot generalize across all three tasks simultaneously. Inverting the cross-attention direction (Q=sMRI, K/V=fMRI) most notably harms MCI vs.\ AD ($0.651{\to}0.515$), validating that fMRI functional context should drive the structural query.

\textbf{Ensemble unimodal comparison.} To assess whether ensembling alone accounts for PCAG's advantage, we additionally trained fMRI-only models with the same five random seeds and evaluated using the same aggregation protocol. The fMRI-only ensemble achieves 0.777, 0.675, and 0.500 on NC vs.\ AD, NC vs.\ MCI, and MCI vs.\ AD, respectively (Table~\ref{tab:ablation}). PCAG-ComBat ensemble outperforms across all three tasks. Most strikingly, the fMRI-only ensemble degrades to chance level on MCI vs.\ AD (AUC~$=0.500$), while PCAG achieves 0.672 ($+0.172$). This demonstrates that sMRI structural information, integrated via cross-attention, is essential for discriminating MCI from AD---a clinically critical boundary where functional connectivity alone is insufficient.

\subsection{Site Confound Analysis}
\label{sec:site}

The most compelling evidence of site confounding appears in NC vs.\ MCI: without harmonization, ADNI within-site AUC is 0.222 ($n{=}11$)---well below chance---revealing that the model classifies ADNI subjects as NC not by detecting cognitive impairment, but by recognizing the scanner signature. After label-free ComBat, ADNI within-site AUC for NC vs.\ MCI rises to 0.444 ($\Delta{=}{+}0.222$), confirming that harmonization restores genuine above-chance discrimination on held-out ADNI subjects (Fig.~\ref{fig:siteconfound}). For NC vs.\ AD, the small ADNI AD test cohort ($n{\approx}3$) precludes reliable within-site interpretation---AUC values derived from three subjects carry extremely wide variance---so this task is not the primary evidence for site confound removal. Post-harmonization GAT attention analysis provides a complementary view: four of nine resting-state networks exhibit class-dominant signal over site artifacts, with DMN, SMN, SN, and CereN showing class effect exceeding site effect by 1.1--6.8 percentage points. Residual site bias in VN and FPN reflects the covariance component unaddressed by mean-and-variance ComBat.

\subsection{Report Quality Evaluation}

With ROI-informed reports, RAG significantly improves clinical relevance (Qwen3: $\Delta{=}{+}0.333$, $p{=}0.007$). No significant coherence penalty was observed for either judge ($p{>}0.17$). Gemma3 exhibits score saturation on clinical relevance ($\Delta{=}0$) once reports contain explicit ROI evidence---a rubric-anchoring effect where both conditions surpass Gemma3's implicit threshold; Qwen3's stronger cross-referencing capability provides a more sensitive signal in this regime. The low inter-judge agreement ($r{=}0.07$--$0.27$) is expected: medical report quality is inherently subjective, and multi-judge evaluation is used here for directional consistency rather than score equivalence \cite{ji2023hallucination}. Judge validity was confirmed by classification-consistency analysis: correctly classified patients scored significantly higher on factual accuracy ($p{=}0.014$) and clinical relevance ($p{=}0.013$), demonstrating that scores reflect genuine clinical accuracy rather than text verbosity. Objectively, RAG reports referenced brain-specific ROIs in 91.8\% of cases and included supporting literature citations in 18.4\%, versus 0\% in no-RAG reports.

% ── Case Study Figure ────────────────────────────────────────────────────
\begin{figure*}[t]
\centering
\begin{tcolorbox}[
  enhanced,
  title={\textbf{Case Study: End-to-End Pipeline Output} \hfill
         \small\textit{Patient sub-023\_S\_6334} $\cdot$
         78-year-old male $\cdot$ 14 years education $\cdot$
         \textbf{Clinical label: MCI}},
  colframe=black!70, colback=white, coltitle=white, colbacktitle=black!70,
  fonttitle=\small, top=4pt, bottom=4pt, left=4pt, right=4pt,
  boxrule=0.6pt, arc=3pt]
%% Row 1: Classification + KG
\begin{minipage}[t]{0.48\textwidth}
\begin{tcolorbox}[enhanced, colframe=gray!50, colback=gray!5,
  title={\small\textbf{(1) PCAG-ComBat Classification}},
  coltitle=black, top=2pt, bottom=2pt, left=3pt, right=3pt, boxrule=0.4pt, arc=2pt]
\small
\begin{tabular}{@{}lcc@{}}
\toprule
Task & Prob. & Prediction\\\midrule
NC vs.\ AD  & 0.318 & NC\\
NC vs.\ MCI & 0.623 & \textbf{MCI}~$\checkmark$\\
MCI vs.\ AD & 0.115 & MCI\\\bottomrule
\end{tabular}\\[2pt]
\textbf{OVO Verdict:} \colorbox{yellow!40}{\textbf{MCI}} (medium confidence)
\end{tcolorbox}
\end{minipage}\hfill
\begin{minipage}[t]{0.48\textwidth}
\begin{tcolorbox}[enhanced, colframe=gray!50, colback=gray!5,
  title={\small\textbf{(2) Knowledge Graph + RAG Retrieval}},
  coltitle=black, top=2pt, bottom=2pt, left=3pt, right=3pt, boxrule=0.4pt, arc=2pt]
\small
\textbf{Similar patients (Neo4j KG):}\\
\quad sub\_0117 (MCI, sim\,$=$\,0.818),\quad 011\_S\_7048 (NC, sim\,$=$\,0.805)\\[2pt]
\textbf{Retrieved literature (MMR, $k{=}5$):}\\
\quad Vemuri et al.~(2014) — Cognitive reserve \& MCI\\
\quad Zhao et al.~(2022) — fMRI graph network for dementia\\
\quad Jack et al.~(2010) — AD biomarker cascade
\end{tcolorbox}
\end{minipage}

\vspace{4pt}
%% Row 2: Neural Evidence
\begin{tcolorbox}[enhanced, colframe=gray!50, colback=gray!5,
  title={\small\textbf{(3) Multimodal Neural Evidence}},
  coltitle=black, top=2pt, bottom=2pt, left=3pt, right=3pt, boxrule=0.4pt, arc=2pt]
\small
\begin{minipage}[t]{0.47\textwidth}
\textbf{sMRI — BrainIAC ViT Attention Rollout:}\\[1pt]
\begin{tabular}{@{}lr@{}}
ROI & Attn.\\\hline
Cingulum\_Mid\_L & 0.288\\
Cingulum\_Post\_R & 0.155\\
Cingulum\_Post\_L & 0.144\\
Precuneus\_R & 0.142\\
Caudate\_L & 0.070\\
\end{tabular}
\end{minipage}\hfill\vline\hfill
\begin{minipage}[t]{0.47\textwidth}
\textbf{fMRI — GAT Network-Level Attention:}\\[1pt]
\begin{tabular}{@{}lr@{}}
ROI & Salience\\\hline
Cingulum\_Mid\_R & 1.000\\
Rolandic\_Oper\_L & 0.978\\
Temporal\_Sup\_R & 0.973\\
Cingulum\_Mid\_L & 0.961\\
Precuneus\_R & 0.953\\
\end{tabular}
\end{minipage}
\end{tcolorbox}

\vspace{4pt}
%% Row 3: Report excerpt
\begin{tcolorbox}[enhanced, colframe=gray!50, colback=blue!3,
  title={\small\textbf{(4) Gemma3-12B Generated Report (excerpt, translated from Traditional Chinese)}},
  coltitle=black, top=2pt, bottom=2pt, left=3pt, right=3pt, boxrule=0.4pt, arc=2pt]
\small\itshape
``BrainIAC attention concentrates on bilateral Cingulum (mid/post, attn\,$=$\,0.29--0.14) and
Precuneus (0.14), structural regions implicated in early MCI and default-mode network
degeneration. GAT functional analysis reveals concordant disruption in Cingulum and
superior temporal cortex (salience\,$\geq$\,0.95), consistent with known salience-network
vulnerability in MCI \cite{vemuri2014cognitive}. Structural and functional evidence
are mutually reinforcing, supporting the \textbf{MCI} classification.
Given the patient's educational attainment (14 years), preserved cognitive reserve may
partially compensate for early network disruption \cite{vemuri2014cognitive}.
\textbf{Recommendation:} Longitudinal follow-up every 6--12 months; monitor bilateral
Cingulum atrophy rate as a progression biomarker.''
\end{tcolorbox}
\end{tcolorbox}
\caption{End-to-end pipeline case study for a correctly classified MCI patient. The system
integrates (1) PCAG-ComBat OVO classification, (2) KG-retrieved similar patients and
RAG-retrieved literature, (3) multimodal neural evidence from BrainIAC ViT attention
rollout (sMRI) and GAT network-level attention (fMRI), and (4) a fact-grounded
Gemma3-12B diagnostic report. Concordant Cingulum and Precuneus involvement across
both modalities supports MCI staging; the report grounds its recommendation in
retrieved clinical literature.}
\label{fig:casestudy}
\end{figure*}

\section{Discussion}

Label-free ComBat proves essential under site-class imbalance: label-aware harmonization preserves site-correlated variation as apparent biological signal, while label-free removal eliminates the confound. The clearest evidence is NC vs.\ MCI, where ADNI within-site AUC rises from 0.222 (below chance) to 0.444 after harmonization---a $\Delta{=}{+}0.222$ shift that marks the transition from site-identity classification to genuine disease-signal discrimination. The NC vs.\ AD within-site result is not interpretable as site confound evidence due to the extremely small ADNI AD test cohort ($n{\approx}3$); with so few positive cases, AUC variance renders any single value uninformative.

The PCAG--Concat contrast reveals task-dependent fusion dynamics. For NC vs.\ AD, the large inter-group divergence makes global concatenation sufficient; for NC vs.\ MCI and MCI vs.\ AD, per-patient cross-modal gating outperforms concatenation by selectively integrating structural context guided by each patient's functional state. The MCI vs.\ AD gap relative to ADMGCN (0.672 vs.\ 0.758) is attributable to our small AD test cohort ($n{=}18$, CI [0.468, 0.879]); the overlapping confidence intervals preclude definitive conclusions.

For clinical reporting, Gemma3's saturation at a score of 2 on clinical relevance reflects an anchoring effect---once reports contain explicit ROI evidence, Gemma3 does not further discriminate based on literature grounding. Qwen3's stronger cross-referencing capability reveals the genuine benefit of RAG: when reports provide concrete anatomical anchors (e.g., Precuneus, Cingulum Mid), retrieved literature integrates meaningfully, yielding significantly higher clinical relevance ($p{=}0.007$). The absence of a coherence penalty resolves the information-density tradeoff observed in earlier pipeline iterations.

\textbf{Limitations.} (1) The held-out test set ($n{=}49$, AD test $n{=}18$) yields wide confidence intervals; MCI vs.\ AD conclusions should be interpreted cautiously. (2) While CovBat~\cite{chen2022covbat} would additionally harmonize covariance structure, our dataset's extreme site-class imbalance (ADNI 70\% NC, TPMIC 69\% disease-stage) poses a fundamental challenge: the between-class covariance difference is confounded with the between-site covariance difference, so CovBat's PCA-based covariance alignment risks removing class-discriminative variance alongside site-specific variance---the same failure mode as label-aware ComBat but at the covariance level. Label-free ComBat's mean-and-variance intervention is more conservative and interpretable under this regime; CovBat remains as future work for datasets with more balanced site-class distributions. (3) The TPMIC dataset is not publicly available, which limits full external reproducibility. (4) Gemma3-12B was chosen for local deployment and data privacy; comparisons with proprietary LLMs are deferred to future work. (5) Clinical reports are currently generated in Traditional Chinese; adaptation to other languages requires re-prompting only.

\section{Conclusion}

We presented an end-to-end multimodal AI pipeline for AD staging that addresses site harmonization, multimodal fusion, and knowledge-grounded report generation jointly. Label-free ComBat eliminates site-class confounds that standard harmonization preserves; PCAG cross-attention provides patient-specific functional--structural integration with consistent advantages on subtle classification tasks; and a Neo4j KG-RAG system grounds Gemma3-12B diagnostic reports in functional network evidence and retrieved medical literature, significantly improving clinical relevance ($\Delta{=}{+}0.333$, $p{=}0.007$). The system further integrates BrainIAC ViT attention rollout for patient-specific sMRI saliency visualization, providing structural interpretability alongside GAT-derived functional evidence. Future work will explore CovBat for covariance-level harmonization and larger multi-site validation cohorts.

\begin{thebibliography}{00}

\bibitem{kawahara2017brainnetcnn} J. Kawahara \textit{et al.}, ``BrainNetCNN: Convolutional neural networks for brain networks; towards predicting neurodevelopment,'' \textit{NeuroImage}, vol. 146, pp. 1038--1049, 2017.

\bibitem{li2021braingnn} X. Li \textit{et al.}, ``BrainGNN: Interpretable Brain Graph Neural Network for fMRI Analysis,'' \textit{Medical Image Analysis}, vol. 74, p. 102233, 2021.

\bibitem{jiang2020higcn} H. Jiang \textit{et al.}, ``Hi-GCN: A hierarchical graph convolution network for graph classification by node and graph-level learning,'' \textit{IEEE J. Biomed. Health Inform.}, vol. 24, no. 12, pp. 3438--3453, 2020.

\bibitem{sun2025admgcn} X. Sun, J. Li, G. Yan, and R. Han, ``ADMGCN: Graph convolutional network for Alzheimer's disease diagnosis with a meta-learning paradigm,'' \textit{Bioinformatics}, 2025, doi: 10.1093/bioinformatics/btaf580.

\bibitem{wen2025stdfbn} S. Wen, J. Wang, W. Liu, X. Meng, and Z. Jiao, ``Spatio-temporal dynamic functional brain network for mild cognitive impairment analysis,'' \textit{Frontiers in Neuroscience}, vol. 19, 2025, doi: 10.3389/fnins.2025.1597777.

\bibitem{yuan2026multimodal} J. Yuan, Y. Liu, J. Jin, S. Li, and Y. Zhang, ``Early detection of Alzheimer's disease via multimodal MRI and machine learning,'' \textit{Frontiers in Aging Neuroscience}, vol. 18, 2026, doi: 10.3389/fnagi.2026.1794982.

\bibitem{johnson2007combat} W. E. Johnson, C. Li, and A. Rabinovic, ``Adjusting batch effects in microarray expression data using empirical Bayes methods,'' \textit{Biostatistics}, vol. 8, no. 1, pp. 118--127, 2007.

\bibitem{fortin2018harmonization} J.-P. Fortin \textit{et al.}, ``Harmonization of cortical thickness measurements across scanners and sites,'' \textit{NeuroImage}, vol. 167, pp. 104--120, 2018.

\bibitem{chen2022covbat} A. A. Chen \textit{et al.}, ``Mitigating site effects in covariance for machine learning in neuroimaging data,'' \textit{Human Brain Mapping}, vol. 43, no. 4, pp. 1179--1195, 2022.

\bibitem{selvaraju2017grad} R. R. Selvaraju \textit{et al.}, ``Grad-CAM: Visual explanations from deep networks via gradient-based localization,'' in \textit{Proc. IEEE/CVF ICCV}, 2017, pp. 618--626.

\bibitem{abnar2020rollout} S. Abnar and W. Zuidema, ``Quantifying attention flow in transformers,'' in \textit{Proc. ACL}, 2020, pp. 4190--4197.

\bibitem{lewis2020rag} P. Lewis \textit{et al.}, ``Retrieval-augmented generation for knowledge-intensive NLP tasks,'' in \textit{Adv. Neural Inf. Process. Syst.}, vol. 33, pp. 9459--9474, 2020.

\bibitem{ji2023hallucination} Z. Ji \textit{et al.}, ``Survey of hallucination in natural language generation,'' \textit{ACM Computing Surveys}, vol. 55, no. 12, pp. 1--38, 2023.

\bibitem{xiong2025medrag} X. Xiong \textit{et al.}, ``MedRAG: Enhancing retrieval-augmented generation with knowledge graph-elicited reasoning for healthcare copilot,'' in \textit{Proc. ACL}, 2025.

\bibitem{wu2025medgraphrag} J. Wu \textit{et al.}, ``Medical Graph RAG: Evidence-based medical large language model via graph retrieval-augmented generation,'' in \textit{Proc. ACL}, 2025.

\bibitem{labkg2025} R. Guo, B. Devereux, G. Farnan, and N. McLaughlin, ``LAB-KG: A retrieval-augmented generation method with knowledge graphs for medical lab test interpretation,'' in \textit{Proc. NeuSymBridge Workshop, COLING}, 2025.

\bibitem{petersen2010adni} R. C. Petersen \textit{et al.}, ``Alzheimer's Disease Neuroimaging Initiative (ADNI): Clinical characterization,'' \textit{Neurology}, vol. 74, no. 3, pp. 201--209, 2010.

\bibitem{tzourio2002aal} N. Tzourio-Mazoyer \textit{et al.}, ``Automated anatomical labeling of activations in SPM using a macroscopic anatomical parcellation of the MNI MRI single-subject brain,'' \textit{NeuroImage}, vol. 15, no. 1, pp. 273--289, 2002.

\bibitem{velickovic2018gat} P. Veli\v{c}kovi\'{c} \textit{et al.}, ``Graph attention networks,'' in \textit{Proc. ICLR}, 2018.

\bibitem{vaswani2017attention} A. Vaswani \textit{et al.}, ``Attention is all you need,'' in \textit{Adv. Neural Inf. Process. Syst.}, vol. 30, 2017.

\bibitem{zhang2024pcag} D. Zhang, R. Nayak, and M. A. Bashar, ``Pre-gating and contextual attention gate: A new fusion method for multi-modal data tasks,'' \textit{Neural Networks}, vol. 179, p. 106553, 2024.

\bibitem{tak2026brainiac} D. Tak \textit{et al.}, ``A generalizable foundation model for analysis of human brain MRI,'' \textit{Nature Neuroscience}, vol. 29, pp. 945--956, 2026.

\bibitem{caruana2004ensemble} R. Caruana, A. Niculescu-Mizil, G. Crew, and A. Ksikes, ``Ensemble selection from libraries of models,'' in \textit{Proc. ICML}, 2004, p. 18.

\bibitem{vemuri2014cognitive} P. Vemuri \textit{et al.}, ``Intellectual and physical activity, personality, and risk of dementia,'' \textit{JAMA Neurology}, vol. 71, no. 11, pp. 1391--1401, 2014.

\bibitem{gemma3_2025} Gemma Team, ``Gemma 3 Technical Report,'' \textit{arXiv preprint arXiv:2503.19786}, 2025.

\end{thebibliography}

\end{document}
